\documentclass[11pt]{article}

\usepackage[margin=1in]{geometry}
\usepackage[T1]{fontenc}
\usepackage{lmodern}
\usepackage{microtype}
\usepackage{setspace}
\usepackage{amsmath,amssymb}
\usepackage{booktabs}
\usepackage{tabularx}
\usepackage{enumitem}
\usepackage{graphicx}
\usepackage{xcolor}
\usepackage{caption}
\usepackage{hyperref}
\usepackage{fancyhdr}

\definecolor{navy}{HTML}{12345B}
\hypersetup{
    colorlinks=true,
    linkcolor=navy,
    citecolor=navy,
    urlcolor=navy,
    pdftitle={Stop or Continue? Forecasting the Same Molecular-Dynamics Trajectory},
    pdfauthor={Avelyn Jing}
}

\pagestyle{fancy}
\fancyhf{}
\fancyhead[L]{YAU 2026 Competition Report}
\fancyhead[R]{Same-Trajectory Forecasting}
\fancyfoot[C]{\thepage}
\setlength{\headheight}{14pt}

\onehalfspacing
\setlist[itemize]{topsep=4pt,itemsep=3pt,parsep=0pt}
\setlist[enumerate]{topsep=4pt,itemsep=3pt,parsep=0pt}
\captionsetup{font=small,labelfont=bf}


\title{
\textbf{Stop or Continue?}\\[0.4em]
\Large Forecasting the Later Behavior of the Same Molecular-Dynamics Trajectory
}
\author{Avelyn Jing\\The Harker School}
\date{YAU High School Science Award 2026}

\begin{document}

\hypersetup{pageanchor=false}
\begin{titlepage}
    \centering
    \vspace*{1.3cm}
    {\Large YAU High School Science Award 2026\par}
    \vspace{0.25cm}
    {\large Chemistry -- Computational Physical Chemistry\par}
    \vspace{1.6cm}
    {\huge\bfseries Stop or Continue?\par}
    \vspace{0.5cm}
    {\LARGE\bfseries Forecasting the Later Behavior of the Same Molecular-Dynamics Trajectory\par}
    \vspace{1.4cm}
    {\Large Avelyn Jing\par}
    \vspace{0.2cm}
    {\large The Harker School\par}
    \vfill
    \begin{minipage}{0.88\textwidth}
    \small
    \textbf{Project scope.}
    This project studies how early measurements can help decide whether a running molecular-dynamics simulation should stop or continue. The test systems contain milbemycin in predicted pockets of AlphaFold-derived ABC-transporter models with water and ions, but no membrane. The work evaluates simulation measurement and resource allocation. It does not claim experimental binding, transporter inhibition, antifungal effectiveness, or reversal of drug resistance.
    \end{minipage}
    \vfill
    {\large Competition submission draft\par}
\end{titlepage}
\hypersetup{pageanchor=true}

\begin{abstract}
A single, correctly defined geometric measurement -- not a machine-learned model -- carries most of the predictive power reported in this study. Once ligand motion is measured protein-relative and periodic-boundary-corrected, that one number, read within the first few nanoseconds of a running trajectory, discriminates the trajectory's own later pose-retention outcome about as well as any fitted classifier tested here. The main signal is carried by a transparent geometric observable rather than by model complexity.

Molecular-dynamics simulations are often run for a fixed length even when some trajectories show early pose divergence. This project asks whether early motion can forecast the later behavior of the \emph{same simulation that is already running}.

The current raw-data inventory contains 81 ligand--pocket trajectories spanning 16 proteins. Seventy-three trajectories reached the nominal 100 ns endpoint and have complete 70--100 ns outcome coverage and all required early measurements. The results reported below retain the previously locked 52-trajectory analysis cohort pending full reconciliation of the expanded cohort with every control analysis. At checkpoints from 5 to 30 ns, ligand motion was measured relative to the protein. The locked late outcome was pose retention: median corrected ligand RMSD below 3 \AA{} during $(70,100]$ ns. The primary 20 ns forecast used mean corrected ligand RMSD and corrected ligand-centroid displacement over $(18,20]$ ns.

The transparent corrected-RMSD signal over $(18,20]$ ns discriminated late same-trajectory retention with AUROC 0.860 (95\% protein-clustered CI 0.747--0.951). Importantly, the same measurement over $(3,5]$ ns already achieved AUROC 0.856 (0.742--0.935). At its exploratory protein-held-out 0.5 operating point, stopping 26 of 52 runs at 5 ns would save 47.5\% of nominal simulation time while incorrectly stopping 1 of 13 late-retained trajectories. The corresponding raw 5 ns and 20 ns AUROCs in the expanded 73-trajectory master cohort were 0.835 and 0.825, respectively. Thus, the current data provide no evidence that waiting until 20 ns improves ranking over a 5 ns checkpoint.

A two-observable leave-one-protein-out logistic model at the inherited 20 ns checkpoint had out-of-fold score AUROC 0.838 (0.706--0.942) and balanced accuracy 0.756 (0.568--0.852) at a fixed 0.5 threshold. PBC-naive and ligand-self-aligned controls gave AUROCs 0.550 and 0.507. The corrected signal transferred poorly to a separately launched trajectory (AUROC 0.546). The 5 ns comparison was selected after examining multiple checkpoints and is therefore an exploratory cost-saving estimate, not prospective validation.

The main contribution is not the coordinate-correction procedure itself, which is established practice, but what it reveals once applied: a same-trajectory pose-persistence signal is already detectable by 5 ns, the earliest checkpoint tested, and does not improve monotonically at later checkpoints. That signal transfers strongly within the same running trajectory but only near chance to a separately launched trajectory from the identical docked structure -- a pattern consistent with trajectory-conditioned, path-specific persistence rather than a proven general law of path-dependence. That mechanistic finding is converted into an explicit, bounded compute/false-stop triage policy. It does not establish binding, biological activity, or replica-level molecular stability. The stopping result is a retrospective development-cohort estimate that requires a frozen 5 ns prospective shadow test. A lightweight software tool is planned that applies this measurement directly to a running simulation's saved trajectory frames alongside the live job, reporting a single stop/continue score together with a confidence range, rather than an opaque model prediction.
\end{abstract}

\noindent\textbf{Keywords:} molecular dynamics; early stopping; ligand pose; periodic boundaries; trajectory forecasting

\tableofcontents
\newpage

\section{Motivation and the Specific Problem}
\label{sec:motivation}

ATP-binding-cassette (ABC) transporters are a clinically important fungal drug-resistance mechanism: overexpressed efflux pumps export azole and other antifungal compounds out of the cell, lowering drug exposure at the intended target and undermining first-line therapy in pathogens on the WHO fungal priority list \cite{WHO2022FPPL,li_azole}. Evaluating candidate compounds against predicted-structure ABC-transporter pockets by molecular dynamics is one way to ask whether a docked ligand stays associated with a resistance-conferring pocket long enough to matter, but explicit-solvent MD is expensive, and a multi-pocket, multi-protein screen forces a concrete computational choice: how long must each candidate run before its outcome is distinguishable from noise?

Molecular docking proposes possible ligand poses, and molecular dynamics tests how those poses evolve over time \cite{meng2011docking,hollingsworth2018md}. In a large screen, running every simulation for the same duration can waste substantial computing time on candidates whose outcome is already legible far earlier.

Whether a system's early configuration predicts its long-time outcome is a specific, long-studied question in physical chemistry, not one this project invents. The committor (splitting probability) of a configuration is the probability that a trajectory launched from it reaches a target basin before returning to the starting basin, and it is the theoretically ideal reaction coordinate for a rare-event process \cite{bolhuis2002tps,e2006towards,peters2016reaction}. Estimating a committor rigorously requires many independent shooting trajectories launched from each configuration of interest -- exactly the repeated-launch cost a screening campaign cannot afford. The practical question studied here is a cheaper, narrower proxy for that idea:

\begin{quote}
\textbf{Based on the trajectory observed so far, should this same simulation stop or continue?}
\end{quote}

This is different from asking whether one short simulation predicts a separately launched simulation of the same complex. The latter probes whether the outcome is a property of the starting structure rather than of the specific stochastic realization already running, though it is not itself a formal committor estimate (Section~\ref{sec:results} states precisely why). For the narrow candidate-triage task studied here, the relevant target is the continuation of the run currently being monitored, and the Results below test both questions on the identical cohort and find different answers: same-trajectory persistence is strong, while transfer to a separately launched trajectory is not. Separate launches remain necessary when the scientific question concerns reproducibility of the molecular complex rather than monitoring of one running job.

\section{Prior Work and This Project's Contribution}

Table~\ref{tab:prior-art} separates established tools, this author's own prior early-termination work on the same screening system (Section~\ref{sec:companion}), and what is actually new in this report, so the specific, original contribution is stated plainly rather than left implicit.

\begin{table}[htbp]
\centering
\caption{What is established, what extends the author's own prior early-termination work (Section~\ref{sec:companion}), and what is new here. Coordinate correction and validation machinery are standard; same-run early-termination triage itself was already introduced in the author's prior 20 ns work; the new contribution is the 100 ns cross-timescale extension, the physical-coordinate ablation as a central result, the model-complexity audit, and the same-run-vs-separate-launch boundary test.}
\label{tab:prior-art}
\small
\begin{tabularx}{\textwidth}{Xl}
\toprule
\textbf{Component} & \textbf{Status} \\
\midrule
Kabsch alignment \cite{kabsch1976} & Established \\
Minimum-image/PBC reconstruction & Established, standard MD practice \\
RMSD and centroid displacement as closed-form geometric observables & Established \\
Logistic regression; leave-one-protein-out validation & Established \\
Transition-path sampling and committor theory \cite{bolhuis2002tps,e2006towards,peters2016reaction} & Established; motivates but is not estimated by this study \\
Same-run early-termination triage on this transporter panel & Extends the author's own prior 20 ns work \cite{jing2026embs} (Section~\ref{sec:companion}) \\
Coordinate-correction ablation as a central, dedicated result & New in this report \\
100 ns cross-timescale extension of the same-run persistence claim & New in this report \\
Same-run vs.\ separate-launch boundary test & New in this report \\
Model-complexity audit (raw observable vs.\ fitted models) & New in this report \\
Compute/false-stop frontier at multiple checkpoints & Extends the same policy structure as prior work, on the new 100 ns data \\
\bottomrule
\end{tabularx}
\end{table}

\section{Relationship to Companion Early-Termination Work}
\label{sec:companion}

This author has two prior manuscripts on early termination of molecular dynamics for the same antifungal-resistance-associated ABC-transporter screening problem, using an earlier 20 ns simulation batch: an IEEE EMBS 2026 submission introducing a physics-informed early-termination framework using pocket RMSD drift/variance and C$\alpha$-displacement descriptors across 72 pockets from 16 transporters \cite{jing2026embs}, and an ACM-BCB 2026 submission that narrows that framework to an FPR-budgeted terminate/continue policy on 80 milbemycin pocket simulations from the same 16 transporters, with the strongest decision window at 2--7 ns. \textbf{The exact publication status of both (published, accepted/forthcoming, submitted, or under review) must be confirmed and stated accurately before submission; it is not asserted here.}

Table~\ref{tab:companion} states the comparison directly; Table~\ref{tab:prior-art} above already attributes which specific components are established, extend prior work, or are new here, and is the authoritative ledger for that distinction.

\begin{table}[htbp]
\centering
\caption{Relationship to this author's prior 20 ns early-termination manuscripts on the same transporter panel. Descriptor and window details for the prior work are drawn directly from their abstracts.}
\label{tab:companion}
\small
\setlength{\tabcolsep}{4pt}
\begin{tabular}{p{0.20\textwidth} p{0.37\textwidth} p{0.37\textwidth}}
\toprule
\textbf{Aspect} & \textbf{Prior 20 ns work \cite{jing2026embs}} & \textbf{This 100 ns report} \\
\midrule
Simulation batch & 20 ns batch & Separate 100 ns batch \\
Prediction target & Same-run 20 ns operational endpoint & Same-run 70--100 ns pose retention \\
Prefix / checkpoint & 1--14 ns (2--7 ns strongest) & 5--30 ns \\
Descriptors & Pocket RMSD drift/variance, C$\alpha$ step statistics & Protein-C$\alpha$-aligned, PBC-corrected ligand pose RMSD and centroid displacement \\
Coordinate ablation & Not a primary result & Central result (Section~\ref{sec:features}) \\
Separate-launch transfer & Not tested & Tested (this section) \\
Model-complexity audit & Not performed in this form & Added (raw observable vs.\ fitted models) \\
Main new conclusion here & Short-run FPR-budgeted triage & Cross-timescale persistence, coordinate dependence, and a tested transfer boundary \\
\bottomrule
\end{tabular}
\end{table}

\section{Main Contribution}

Building on Table~\ref{tab:prior-art}, the project makes three linked contributions:

\begin{enumerate}
    \item It repurposes a closed-form, physically defined pose observable -- not something learned from data -- as an early-commitment triage signal read directly from a currently running trajectory, rather than as a post-hoc pose-validation metric.
    \item It provides direct empirical evidence, via a same-trajectory-vs-separate-launch comparison on the identical cohort, that pose commitment here is \emph{path-dependent} -- a property of the specific stochastic realization already running -- rather than fixed by the docked starting structure alone.
    \item It converts that mechanistic finding into an explicit compute/false-stop frontier: because commitment cannot be read from the structure or from a second launch, only the running trajectory's own early dynamics can support triage, and the resulting compute saving is reported together with its false-stop cost.
\end{enumerate}

Figure~\ref{fig:core-idea} shows the temporal design: an early observation is separated from the late outcome by a forecast gap.

\begin{figure}[htbp]
    \centering
    \includegraphics[width=0.96\textwidth]{figures/core_idea.png}
    \caption{Core same-trajectory forecasting idea. The same corrected ligand RMSD is measured twice in one running trajectory: an early window, $(18,20]$ ns, and a late window, $(70,100]$ ns, separated by a forecast gap. The locked outcome is binary -- late pose retained if median corrected RMSD over the late window stays below 3~\AA, late pose non-retained otherwise -- not an experimental binding or stability label.}
    \label{fig:core-idea}
\end{figure}

The central scientific result is:

\begin{quote}
\textbf{In this retrospective cohort, a protein-relative, periodic-boundary-consistent ligand-pose signal was already detectable by the 5 ns checkpoint and discriminated late pose retention within the same running trajectory. The same descriptor did not transfer reliably from a separately launched trajectory, indicating that the forecast is trajectory-conditioned rather than an intrinsic stability prediction for the starting complex.}
\end{quote}

\section{Data and Forecasting Design}

The systems contained milbemycin docked into predicted pockets of AlphaFold-derived ABC-transporter structures. Pockets were identified using P2Rank, and docking poses were generated using AutoDock Vina \cite{jumper2021alphafold,mirdita2022colabfold,krivak2018p2rank,trott2010vina}.

\subsection{Molecular-dynamics simulation protocol}

The docked complexes were prepared as explicit-solvent OpenMM \cite{eastman2024openmm} systems. Protein atoms used AMBER14/\texttt{ff14SB} \cite{maier2015ff14sb}; milbemycin used GAFF-2.11 parameters \cite{wang2004gaff} generated through \texttt{openmmforcefields}, with AM1--BCC partial charges \cite{jakalian2002am1bcc} assigned by the OpenFF Toolkit. Each system used TIP3P water \cite{jorgensen1983tip3p} in a rectangular box with at least 1.0~nm protein--box padding and 0.15~M NaCl ionic strength through \texttt{Modeller.addSolvent}. No lipid bilayer was included. Long-range electrostatics used PME with a 1.0~nm real-space cutoff; the Lennard--Jones interaction used a 0.8~nm switching distance with dispersion correction. Bonds to hydrogen were constrained (\texttt{HBonds}), hydrogen masses were repartitioned to 4~Da, and the LangevinMiddle integrator used a 4~fs timestep at 300~K with friction 1~ps$^{-1}$. A Monte Carlo barostat maintained 1~atm and attempted volume moves every 25 steps.

For each nominal 100~ns run, the system was energy-minimized for up to 1,000 iterations, equilibrated for 25,000 steps (100~ps), and then propagated for 25,000,000 production steps. The DCD trajectory was written every 50,000 steps (200~ps), giving 500 nominal frames per 100~ns; the online fingerprint reporter additionally sampled every 5,000 steps (20~ps) and emitted 2~ns summaries. The saved cohort contains one 100~ns \texttt{replica\_1} launch per complex. The launch code requested a CUDA GPU, mixed precision, and deterministic forces. The separate-launch control was a separately solvated, minimized, and equilibrated 20~ns launch from the same prepared complex, not a second 100~ns replicate. The repository requirements snapshot pins OpenMM 8.2.0, but no per-run environment manifest links that pin to every saved trajectory; production logs do not record the runtime package version or GPU model. Per-launch random seeds were also not durably preserved, so paired runs are described as separately launched rather than independently seeded. Seven trajectories ended before the nominal 100~ns window and were excluded from the locked complete-coverage cohort, so the 500-frame count is a target protocol rather than a guarantee for every saved run.


\begin{table}[t]
  \centering
  \caption{Screening protocol and endpoint definition for the early MD triage study.}
  \label{tab:study_design_summary}
  \small
  \setlength{\tabcolsep}{4pt}
  \renewcommand{\arraystretch}{1.0}
  \begin{tabular}{p{0.28\columnwidth} p{0.68\columnwidth}}
    \toprule
    \textbf{Design choice} & \textbf{Setting and purpose} \\
    \midrule

    Production endpoint &
    Nominal 100\,ns trajectory; complete coverage is required for the late outcome. \\

    Forecast checkpoints &
    5, 10, 15, 20, and 30\,ns, each using corrected RMSD over the preceding 2\,ns. The 5\,ns checkpoint is the earliest evaluated. \\

    Late outcome &
    Pose retention requires complete $(70,100]$\,ns coverage and median protein-relative, PBC-corrected ligand RMSD below 3~\AA. \\

    Prediction scope &
    Continuation of the same running trajectory; not binding, biological activity, equilibrium stability, or replica transfer. \\

    Primary split &
    Leave-one-protein-out validation tests transfer to unseen proteins. \\

    Decision threshold &
    Fixed model score of 0.5; uncalibrated and not selected to satisfy a false-stop constraint. \\

    Reported risk &
    Protein-clustered AUROC intervals, out-of-fold decision counts, false stops, and nominal compute saving. \\

    Validation status &
    Retrospective checkpoint comparison; a frozen 5\,ns prospective shadow test is required before deployment. \\

    \bottomrule
  \end{tabular}
\end{table}


The audit began with 64 discovered systems. Sixty-one had usable separate-launch 20 ns measurements, 59 had cached late metrics, and 52 had complete nominal $(70,100]$ ns coverage. Those 52 trajectories, spanning 15 proteins with 13 late-retained and 39 late-non-retained outcomes, form the locked cohort used for all headline comparisons. Early measurements were evaluated at 5, 10, 15, 20, and 30 ns. The outcome came from the final 30 ns:

\[
\text{early checkpoint }t_c < 70\ \mathrm{ns}
\quad\longrightarrow\quad
\text{late outcome from }70\text{--}100\ \mathrm{ns}.
\]

No frame from $(70,100]$ ns was used to construct an early measurement. The 20 ns decision time was inherited from the pre-existing screening workflow; it was not selected by maximizing the five-checkpoint curve. The 5, 10, 15, and 30 ns analyses are therefore secondary retrospective explorations. This report freezes 20 ns for a future prospective shadow test, but does not claim that it is an optimal checkpoint.

The late label describes geometric pose retention:

\begin{equation}
Y_L = \mathbf{1}\!\left[
\text{complete }(70,100]\text{ ns coverage}
\;\land\;
\operatorname*{median}_{t\in(70,100]}d_{\mathrm{RMSD}}(t)<3\ \text{\AA}
\right].
\end{equation}

The earlier two-condition label also required centroid displacement below 5 \AA. That condition is redundant: for these atom-displacement definitions, ligand-centroid displacement cannot exceed ligand RMSD. Across all 59 cached late metrics, the audit found no inequality violations and no label changes after removing the centroid condition. Centroid displacement remains a continuous early predictor and diagnostic, but not an independent late-label condition. The outcome is a simulation geometry label, not an experimental binding or activity label.

\section{Physical Measurement and Its Meaning}
\label{sec:features}

The coordinate-correction procedure below is established MD practice (Table~\ref{tab:prior-art}: Kabsch alignment, minimum-image PBC reconstruction). What this section and the ablations in Section~\ref{sec:results} test is not whether this procedure exists, but whether skipping it changes the outcome -- and, more importantly, whether the resulting corrected descriptor reveals path-dependent commitment once applied correctly.

Figure~\ref{fig:primary-features} summarizes the coordinate processing and the two primary measurements. Pocket-neighborhood retention supports structural interpretation. Early residual-motion trend is a separate exploratory warning signal and is not part of the locked primary model.

\begin{figure}[p]
    \centering
    \includegraphics[width=0.94\textwidth]{figures/primary_features.png}
    \caption{Primary observable calculation. Every frame is made PBC-consistent and aligned to frame 0 using protein C$\alpha$ atoms; the protein-derived alignment transform is then applied to ligand heavy atoms. Corrected ligand RMSD and the unweighted ligand-centroid displacement are averaged over $(18,20]$ ns. The locked late label additionally requires complete nominal coverage over $(70,100]$ ns and median corrected RMSD below 3~\AA. The diagram uses ``probability'' for the logistic output; because the class-weighted model was not calibrated, this report interprets that output as a model score rather than an absolute probability.}
    \label{fig:primary-features}
\end{figure}

\subsection{Correcting the trajectory}

Two coordinate effects can create false ligand motion. First, the protein and ligand can rotate or translate together. Second, periodic-boundary wrapping can make a ligand appear to jump across the box.

For every frame, protein C$\alpha$ atoms were centered and Kabsch-aligned to frame 0 \cite{kabsch1976}. Ligand heavy atoms were reconstructed in the periodic image nearest the protein using the minimum-image convention. The protein-derived alignment transform was then applied to the ligand. The ligand was \emph{not} aligned to itself, because self-alignment would remove its movement relative to the pocket.

\subsection{The corrected-pose observable}

Corrected ligand RMSD measures atomic-pose change:

\begin{equation}
d_{\mathrm{RMSD}}(t)=
\sqrt{\frac{1}{N_L}\sum_{a=1}^{N_L}
\left\|\widetilde{\mathbf r}_a(t)-\widetilde{\mathbf r}_a(0)\right\|^2}.
\end{equation}

Corrected ligand-centroid displacement measures translation:

\begin{equation}
d_{\mathrm{disp}}(t)=
\left\|\mathbf C_L(t)-\mathbf C_L(0)\right\|_2.
\end{equation}

Here, $\mathbf C_L$ is the unweighted centroid of the selected ligand heavy atoms. $d_{\mathrm{RMSD}}(t)$ and $d_{\mathrm{disp}}(t)$ are not learned from data: they are closed-form functions of the trajectory coordinates with no fitted or tunable parameters. Kabsch alignment has an exact least-squares solution via singular value decomposition \cite{kabsch1976}; the minimum-image convention is a fixed geometric rule for periodic systems; and the two equations above are direct geometric measurements on the corrected coordinates. Nothing about these quantities is learned from labels, and they are computed identically whether or not any classifier is ever fitted on top of them. In this sense they function as a physically defined observable -- close in spirit to a reaction-coordinate proxy for the committor introduced in Section~\ref{sec:results} -- rather than as engineered input to a model.

The only learned component in this project is a two-parameter logistic combination of the standardized observable, used solely to produce the retrospective stopping policy in Section~\ref{sec:results}. At 20 ns, both quantities are averaged over $(18,20]$ ns before that combination is fit. Evaluation holds out all complexes from one protein at a time. Within each outer fold, a standard scaler is fitted only on the training proteins and applied unchanged to the held-out protein. The fixed estimator is L2 logistic regression with $C=1.0$, balanced class weights, an intercept, and the \texttt{lbfgs} solver. No hyperparameter tuning or inner cross-validation was performed, and Section~\ref{sec:results} shows this fitted combination does not exceed the raw observable's own discriminative power -- the scientific content is carried by the equation, not by the classifier.

The continue decision uses a fixed score threshold of 0.5. This threshold was not selected to satisfy a false-stop constraint. Every reported policy decision is out of fold, but the observable-set and alternative-checkpoint comparisons were examined retrospectively on the development cohort. Leave-one-protein-out evaluation therefore yields 15 fold-specific evaluation models; it does not produce a final deployment model trained on all development proteins.

\subsection{Post-lock model-complexity audit}

After the primary analysis was locked, a separate exploratory model-complexity audit was performed on the expanded complete-coverage inventory of 73 trajectories from 16 proteins. All candidate inputs ended at 20 ns. Of the 21 trajectories added beyond the locked 52, 20 are additional pockets or replicas on proteins already represented in the locked cohort; only one protein (a 16th transporter) is entirely new to this expansion. Class balance held steady across the growth: 13/39 (25.0\%) late-retained at $n=52$ versus 17/56 (23.3\%) at $n=73$. This audit is therefore best read as added statistical power and a stability check on already-sampled proteins, not as a test of generalization to a substantially broader protein panel; the primary cohort is being expanded further toward 80 trajectories, and the same-trajectory-vs-separate-launch comparison in Section~\ref{sec:results} is not yet re-derived at that larger scale. The audit compared the raw corrected-RMSD score and the fixed two-observable endpoint logistic model with temporal logistic regression, a quadratic-spline logistic model, histogram gradient boosting, and random forest. Temporal models received corrected RMSD and centroid displacement at 5, 10, 15, and 20 ns together with changes and discrete acceleration summaries; no post-20 ns measurement was used as a predictor.

Evaluation used an outer leave-one-protein-out loop. For every model requiring tuning, hyperparameters were selected using only the outer training proteins through an inner protein-grouped leave-one-protein-out comparison. The held-out protein therefore contributed neither to model fitting nor to configuration selection. This audit was conducted after inspection of the locked results and is reported as exploratory model auditing rather than prospective validation.

\subsection{Supporting pocket-neighborhood retention}

A baseline pocket neighborhood $A_0$ is formed from residues that contact the ligand in at least 20\% of frames during 0--1 ns, using a 4.5 \AA{} heavy-atom cutoff. At each frame,

\begin{equation}
R(t)=\frac{1}{|A_0|}\sum_{i\in A_0}\mathbf{1}[d_i(t)\leq4.5\ \text{\AA}].
\end{equation}

The 20 ns pocket-retention measurement is the mean of $R(t)$ over $(18,20]$ ns. It helps distinguish a ligand that has moved within its original neighborhood from one that has lost that neighborhood. On the locked cohort its same-trajectory AUROC was 0.714 (95\% CI 0.560--0.850). It is reported separately as supporting evidence, not as an input to the locked primary stopping model.

\subsection{Exploratory early residual-motion trend}

Early Residual Drift (ERD) summarizes slopes of aligned pocket RMSD over the first few nanoseconds:

\begin{equation}
\mathrm{ERD}_{3\text{--}5}=\frac{\mathrm{slope}_{3\mathrm{ns}}+\mathrm{slope}_{4\mathrm{ns}}+\mathrm{slope}_{5\mathrm{ns}}}{3}.
\end{equation}

ERD was explored as a later non-retention warning among trajectories without an early pose warning. Because its threshold was selected using the current non-retained cases, it is not included in the locked primary forecast and requires prospective validation.

\section{Results}
\label{sec:results}

Unless otherwise noted, intervals below use 5,000 protein-clustered bootstrap draws. They resample whole proteins while treating the continuous or saved out-of-fold scores as fixed; they do not refit the complete modeling and selection procedure within each draw.

Figure~\ref{fig:locked-results} consolidates the latest locked findings. Taken together, they show that corrected early pose is a strong same-run persistence signal, coordinate-inconsistent measurements lose most of that discrimination, one separately launched trajectory does not provide comparable predictive transfer, and any compute saving must be reported together with its false stops.

\begin{figure}[p]
    \centering
    \includegraphics[width=\textwidth,height=0.72\textheight,keepaspectratio]{figures/latest_locked_results.pdf}
    \caption{Latest locked same-trajectory findings on the identical complete-coverage cohort: 52 trajectories from 15 proteins, including 13 late-retained and 39 late-non-retained outcomes. (A) Mean corrected RMSD over $(18,20]$ ns is associated with median corrected RMSD over $(70,100]$ ns in the same run; the displayed Spearman association is descriptive and is not adjusted for protein clustering. (B) Corrected protein-relative, PBC-aware RMSD discriminates the locked late outcome, whereas PBC-naive and ligand-self-aligned controls are near chance. (C) Corrected pose is informative for the same running trajectory but transfers poorly from one separately launched trajectory; the pocket-retention difference is less conclusive. (D) The retrospective policy frontier makes the compute--error tradeoff explicit: at the fixed 0.5 score threshold, nominal compute saving is 43.1\% with 2 of 13 late-retained trajectories stopped. AUROC intervals use 5,000 protein-clustered bootstrap draws with fixed scores; the complete model and analysis-selection procedure is not refitted. The policy frontier is retrospective, not a prospective guarantee.}
    \label{fig:locked-results}
\end{figure}

\subsection{The corrected-RMSD signal appears at every explored checkpoint}

\begin{table}[htbp]
\centering
\caption{Exploratory checkpoint comparison on the identical locked cohort ($n=52$, 15 proteins). Each row uses mean corrected RMSD over the final 2 ns before the named checkpoint. Balanced accuracy uses the saved protein-held-out single-observable logistic score and fixed 0.5 threshold.}
\label{tab:checkpoint}
\begin{tabular}{lcc}
\toprule
Checkpoint & Corrected-RMSD AUROC (95\% CI) & LOPO balanced accuracy \\
\midrule
5 ns  & 0.856 [0.742, 0.935] & 0.782 \\
10 ns & 0.888 [0.771, 0.971] & 0.859 \\
15 ns & 0.884 [0.764, 0.968] & 0.846 \\
20 ns & 0.860 [0.747, 0.951] & 0.808 \\
30 ns & 0.854 [0.742, 0.940] & 0.821 \\
\bottomrule
\end{tabular}
\end{table}

\begin{figure}[htbp]
    \centering
    \includegraphics[width=0.78\textwidth]{figures/checkpoint_curve.pdf}
    \caption{Corrected-RMSD AUROC by checkpoint, from Table~\ref{tab:checkpoint} ($n=52$, 15 proteins). The confidence intervals overlap; this analysis does not establish an optimal checkpoint. It shows that useful information is present by the earliest tested checkpoint and does not increase monotonically with time.}
    \label{fig:checkpoint-curve}
\end{figure}

The earliest tested checkpoint, $(3,5]$ ns, achieved nearly the same discrimination as $(18,20]$ ns. In the expanded 73-trajectory master cohort, the corresponding descriptive raw AUROCs were 0.835 at 5 ns and 0.825 at 20 ns, so the early result was not confined to the original 52 rows. The 10 ns point estimate was numerically highest in the locked cohort, but the intervals overlap substantially and checkpoint selection was not nested. The study therefore does not identify an optimal checkpoint. Twenty nanoseconds remains the inherited primary checkpoint; 5 ns is the strongest candidate to freeze for the next prospective test.

This finding changes the practical interpretation. The physical persistence signal need not require observing one fifth of a 100 ns run: in this dataset it was already measurable after approximately 5 ns. The result supports an earlier triage hypothesis, not a claim that every ligand has equilibrated by 5 ns.

\subsection{The physical signal and learned decision score are distinct}

Across the locked cohort, early and late corrected RMSD were strongly associated (descriptive Spearman $\rho=0.73$), consistent with temporal persistence of the same geometric observable. This row-level correlation was not adjusted for protein clustering.

At 20 ns, the transparent corrected-RMSD signal had AUROC 0.860 (95\% CI 0.747--0.951). Using only the saved frame nearest 20 ns gave the same AUROC point estimate, 0.860 (0.738--0.947), indicating that the result was not created by averaging over the preceding 2 ns; this single-frame control remains a distinct estimand. The two-observable logistic model combined corrected RMSD with corrected centroid displacement. Its saved outer-LOPO score had AUROC 0.838 (0.706--0.942). At the fixed 0.5 threshold, balanced accuracy was 0.756 (0.568--0.852), sensitivity to late retention was 0.846, and specificity for late non-retention was 0.667.

The raw corrected-RMSD AUROC and the two-observable model-score AUROC answer different questions. The first measures discrimination by one physically defined observable; the second evaluates a fitted decision mapping. The raw RMSD point estimate was higher, so this report does not claim that logistic regression itself adds scientific value. The main result is the coordinate-consistent measurement and its same-run scope.

\subsection{The equation matches or beats every learned model tested}

The post-lock expanded-cohort audit tested whether nonlinear or temporal learning could recover additional predictive structure beyond the endpoint measurement. Table~\ref{tab:model-audit} and Figure~\ref{fig:physics-vs-ai} show that it did not. The raw corrected-RMSD score -- the closed-form equation defined in Section~\ref{sec:features}, with no fitted parameters -- remained the strongest ranking signal, with AUROC 0.825 (95\% protein-clustered CI 0.733--0.900). Every fitted model tested, including a two-observable logistic regression given the same inputs, scored at or below it: the endpoint logistic model achieved 0.818 (0.723--0.895), while temporal logistic regression, random forest, spline logistic regression, and gradient boosting produced lower point estimates still. The paired temporal-minus-endpoint logistic AUROC difference was $-0.021$; a 1,000-draw protein-clustered bootstrap interval was approximately $[-0.115,0.067]$, providing no evidence that trajectory-history measurements, or any of the added model complexity, improved ranking over the equation alone.

\begin{figure}[htbp]
    \centering
    \includegraphics[width=0.92\textwidth]{figures/physics_vs_ai_model_comparison.pdf}
    \caption{The closed-form physics equation (top) against five learned models given the same or richer inputs, on the identical expanded cohort ($n=73$, 16 proteins). No fitted or learned model exceeds the equation's own discrimination; all comparisons use 95\% protein-clustered bootstrap intervals as in Table~\ref{tab:model-audit}. This is offered as a direct visual counterpart to the table, not a separate analysis.}
    \label{fig:physics-vs-ai}
\end{figure}

\begin{table}[htbp]
\centering
\caption{Exploratory model-complexity audit on the expanded complete-coverage cohort ($n=73$, 16 proteins). Tuned models use nested protein-grouped selection inside outer leave-one-protein-out evaluation.}
\label{tab:model-audit}
\begin{tabularx}{\textwidth}{Xcc}
\toprule
Model or score & AUROC & 95\% clustered CI \\
\midrule
Raw corrected RMSD & \textbf{0.825} & [0.733, 0.900] \\
Endpoint logistic regression & 0.818 & [0.723, 0.895] \\
Temporal logistic regression & 0.797 & [0.681, 0.905] \\
Temporal random forest & 0.788 & [0.698, 0.869] \\
Endpoint spline logistic regression & 0.785 & [0.692, 0.868] \\
Temporal histogram gradient boosting & 0.712 & [0.627, 0.820] \\
\bottomrule
\end{tabularx}
\end{table}

At the fixed 0.5 score threshold, temporal logistic regression stopped 42 trajectories, including 3 of 17 late-retained trajectories, whereas endpoint logistic regression stopped 40 with the same 3 false stops. The temporal model therefore stopped two additional late-non-retained trajectories in this retrospective cohort, but its AUROC was lower and its scores were not calibrated. This small policy difference is hypothesis-generating and requires prospective testing.

The model-complexity comparison supports a simple interpretation: once coordinates are made protein-relative and periodic-boundary consistent, the current corrected pose contains most of the recoverable same-run persistence signal. Representation quality mattered more than classifier complexity. Because the predictor and outcome use the same corrected-RMSD family, this finding should be interpreted as low-dimensional temporal persistence, not discovery of a hidden molecular mechanism.

\subsection{Coordinate-consistent processing reveals the signal}

\begin{table}[htbp]
\centering
\caption{Locked 20 ns continuous-score comparisons on the same 52 trajectories.}
\label{tab:locked-comparisons}
\begin{tabularx}{\textwidth}{Xcc}
\toprule
Measurement and prediction task & AUROC & 95\% CI \\
\midrule
Corrected RMSD $\rightarrow$ same trajectory & \textbf{0.860} & [0.747, 0.951] \\
Protein aligned, no PBC correction $\rightarrow$ same trajectory & 0.550 & [0.408, 0.743] \\
Ligand-self-aligned RMSD $\rightarrow$ same trajectory & 0.507 & [0.323, 0.741] \\
Corrected RMSD from a separate launch $\rightarrow$ target outcome & 0.546 & [0.314, 0.674] \\
Pocket retention $\rightarrow$ same trajectory & 0.714 & [0.560, 0.850] \\
Pocket retention from a separate launch $\rightarrow$ target outcome & 0.540 & [0.364, 0.693] \\
\bottomrule
\end{tabularx}
\end{table}

The paired AUROC difference between corrected and PBC-naive RMSD was 0.310 (95\% CI 0.036--0.504); the corrected-minus-ligand-self difference was 0.353 (0.111--0.532). Thus, representations that corrupt periodic motion or discard protein-relative translation and rotation contained substantially less information about the late pose-retention label.

\subsection{The 3 \AA{} outcome threshold is not fragile}

The locked late label requires median corrected RMSD below 3~\AA{} over the last 30\% of the trajectory. Because the early predictor and the late outcome are both drawn from the same corrected-RMSD family, it matters whether this specific cutoff is doing unearned work. Table~\ref{tab:threshold-sensitivity} recomputes the label at five nearby cutoffs, holding the early corrected-RMSD predictor and cohort fixed, and re-evaluates AUROC at both the 5~ns and 20~ns checkpoints.

\begin{table}[htbp]
\centering
\caption{Late-outcome threshold sensitivity on the identical locked cohort ($n=52$, 15 proteins). Only the late-window cutoff defining retention changes across rows; the early corrected-RMSD predictor is unchanged.}
\label{tab:threshold-sensitivity}
\small
\setlength{\tabcolsep}{4pt}
\begin{tabular}{lccc}
\toprule
Late RMSD threshold & Ret./non-ret. & 5 ns AUROC (95\% CI) & 20 ns AUROC (95\% CI) \\
\midrule
2.0 \AA & 3/49  & 0.844 [0.702, 0.964] & 0.816 [0.612, 0.984] \\
2.5 \AA & 6/46  & 0.837 [0.718, 0.944] & 0.877 [0.686, 1.000] \\
3.0 \AA\ (primary) & 13/39 & \textbf{0.856} [0.742, 0.935] & \textbf{0.860} [0.747, 0.951] \\
3.5 \AA & 19/33 & 0.844 [0.722, 0.930] & 0.893 [0.794, 0.973] \\
4.0 \AA & 21/31 & 0.823 [0.643, 0.941] & 0.892 [0.791, 0.974] \\
5.0 \AA & 30/22 & 0.741 [0.581, 0.869] & 0.829 [0.679, 0.947] \\
\bottomrule
\end{tabular}
\end{table}

Every cutoff from 2.0 to 5.0~\AA\ keeps both checkpoint AUROCs well clear of chance, and the point estimates at 20~ns are, if anything, slightly higher at intermediate cutoffs (3.5--4.0~\AA) than at the primary 3.0~\AA\ choice. The class balance shifts substantially across this sweep (from 3 to 30 retained out of 52), and only at the most permissive 5.0~\AA\ cutoff -- where retained outnumber non-retained -- does the 5~ns AUROC noticeably soften. The result does not collapse around the chosen threshold, so the 3.0~\AA\ cutoff is not doing unearned discriminative work; it is a reasonable, non-adversarial choice among several that would have supported the same qualitative claim.

\subsection{Prediction beyond obvious early escape}

A separate concern is whether the same-trajectory signal is trivial in the sense of only separating already-obvious catastrophic early departures from everything else. Table~\ref{tab:non-obvious} re-evaluates the locked 20~ns AUROC on three restricted subsets of the same 52 trajectories, using the unchanged 3~\AA\ locked label.

\begin{table}[htbp]
\centering
\caption{Same-trajectory AUROC at 20 ns on restricted subsets of the locked cohort, using the unchanged 3~\AA\ label. ``Near the decision boundary'' keeps only trajectories with 20 ns corrected RMSD within 1.5~\AA\ of the 3~\AA\ cutoff, i.e.\ the hardest, least obvious cases.}
\label{tab:non-obvious}
\begin{tabular}{lccc}
\toprule
Restricted cohort & $n$ (proteins) & Retained & AUROC (95\% CI) \\
\midrule
Full locked cohort & 52 (15) & 13 & 0.860 [0.747, 0.951] \\
Excluding top 10\% highest early RMSD & 46 (15) & 13 & 0.834 [0.705, 0.939] \\
Middle 80\% of early RMSD (trims both tails) & 40 (15) & 10 & 0.887 [0.750, 0.984] \\
Near the decision boundary ($\pm$1.5~\AA\ of 3~\AA) & 34 (14) & 13 & 0.835 [0.625, 0.957] \\
\bottomrule
\end{tabular}
\end{table}

Removing the most extreme 10\% of early departures leaves the AUROC essentially unchanged (0.834 vs.\ 0.860), and restricting to the middle 80\% -- which removes both the most obvious escapes and the most obviously static poses -- if anything increases it slightly (0.887). The most stringent test keeps only trajectories whose 20~ns corrected RMSD falls within 1.5~\AA\ of the 3~\AA\ cutoff itself, i.e.\ the cases hardest to call by inspection; even there the AUROC remains clearly above chance (0.835, 95\% CI 0.625--0.957, $n=34$). The signal is therefore not an artifact of a few unambiguous early escapes dominating the ranking.

\subsection{The forecast is trajectory-conditioned, not structure-determined}

Corrected RMSD was much more informative for the same running trajectory than for the separately launched trajectory. The paired same-minus-separate AUROC difference was 0.314 (95\% CI 0.181--0.526). In contrast, the pocket-retention difference was 0.174 ($-0.011$--0.386), so that supporting comparison was not conclusive.

Because this comparison depends on the two launches actually being dynamically independent rather than accidental repeats of the same trajectory, that independence was checked directly rather than assumed. For all 61 same-complex pairs, the starting structure files were byte-identical (matching file hashes), so any divergence is not due to a different starting pose. Of these, 54 pairs had sufficient overlapping simulated time to compare potential-energy time series between the two launches; the correlation of energy fluctuations about each trajectory's own trend was indistinguishable from zero (median 0.004, range $-0.074$ to $0.148$, no pair exceeding 0.3 in magnitude). Two identically seeded or otherwise coupled trajectories would show strong energy correlation over their overlapping window; near-zero correlation from an identical starting structure is direct evidence that the two launches are independent stochastic realizations, not a documentation gap papered over with an assumption.

This study is motivated by the broader concept of dynamical commitment studied in transition-path sampling \cite{bolhuis2002tps,e2006towards,peters2016reaction}, but it does not estimate a committor. A formal committor requires an ensemble of trajectories initiated from the \emph{same} checkpoint configuration under a specified momentum or noise distribution; the separate launch here is instead re-solvated, re-minimized, and re-equilibrated from the original $t=0$ docked structure, not shot forward from the 20 ns (or any later) phase-space point of the running trajectory. The operational question actually tested is narrower and cheaper: whether the observed prefix of one already-realized trajectory forecasts its own continuation, versus forecasting an independent realization launched from the same starting structure. The near-chance separate-launch AUROC (0.546) against the strong same-trajectory AUROC (0.860) is consistent with that forecast being trajectory-conditioned -- carried by the specific realization already running -- rather than a property recoverable from the starting heavy-atom coordinates alone; it does not prove path-dependent commitment as a general phenomenon beyond this cohort.

Only one additional launch per complex was examined, so this reading rests on transfer to one \emph{separately launched trajectory} per complex rather than a full replica-level reproducibility distribution, and the energy-decorrelation check covered 54 of 61 pairs rather than all of them. It should be treated as a verified, theory-motivated demonstration for this cohort, not as a general proof.

\subsection{The retrospective stopping point exposes the compute--error tradeoff}

At the fixed 0.5 model-score threshold, the out-of-fold decision table was:

\begin{table}[htbp]
\centering
\caption{Retrospective out-of-fold operating point at 20 ns. ``Continue'' predicts late pose retention.}
\label{tab:policy}
\begin{tabular}{lccc}
\toprule
Decision at 20 ns & Late retained & Late non-retained & Total \\
\midrule
Stop & 2 & 26 & 28 \\
Continue & 11 & 13 & 24 \\
\midrule
Total & 13 & 39 & 52 \\
\bottomrule
\end{tabular}
\end{table}

Stopping 28 trajectories at 20 ns instead of 100 ns would save $28\times80=2{,}240$ of 5,200 nominal simulated ns, or 43.1\%. The policy stopped 26 of 39 late-non-retained trajectories and incorrectly stopped 2 of 13 late-retained trajectories. Thirteen of 39 non-retained trajectories were still continued, so ``continue'' is not a guarantee of later retention.

A budget-matched random baseline makes this concrete: stopping 28 of 52 trajectories at random would be expected to stop $28\times\frac{13}{52}=7.0$ of the 13 late-retained trajectories, versus the 2 actually stopped by the policy at the identical compute budget. At 5 ns, the corresponding random expectation is $26\times\frac{13}{52}=6.5$ against the 1 late-retained trajectory the policy actually stopped. At the same nominal compute budget, the checkpoint policy loses roughly one late-retained trajectory for every three to seven a budget-matched random policy would have lost.

As a check on whether the fitted two-observable model is doing decision-relevant work beyond the raw observable, a direct threshold on corrected RMSD alone -- selected within each leave-one-protein-out fold using only training-protein balanced accuracy, with no logistic fitting at all -- was evaluated at 20 ns. It reproduced the identical confusion matrix (28 stopped: 2 late-retained, 26 late-non-retained; balanced accuracy 0.756; 43.1\% nominal saving). At 5 ns, the same training-fold-optimal direct threshold gave a more aggressive operating point (35 stopped: 3 late-retained, 32 late-non-retained; balanced accuracy 0.795) than the previously reported single-observable logistic score at its fixed 0.5 threshold (26 stopped: 1 late-retained, 25 late-non-retained); both are monotonic in the same underlying corrected-RMSD ranking, so the two differ only in where the operating point falls, not in discrimination. This confirms the fitted logistic combination is not carrying independent scientific content at 20 ns and shows that the specific stop/false-stop tradeoff, at any checkpoint, is a choice of decision rule on top of one observable rather than a property of the observable itself.

This threshold was not selected to enforce a false-stop ceiling, and the score is uncalibrated. These are retrospective development-cohort operating characteristics, not prospective false-stop control or a deployment guarantee.

Across the full retrospective threshold sweep, greater nominal compute saving required accepting a higher false-stop rate; no point on that frontier has yet been selected or confirmed prospectively.

\subsection{An exploratory 5 ns decision could reduce nominal cost further}

The checkpoint comparison also saved out-of-fold predictions from a single-observable logistic model fitted to corrected RMSD at each checkpoint. At the fixed 0.5 threshold, the 5 ns model produced the following retrospective decisions:

\begin{table}[htbp]
\centering
\caption{Exploratory out-of-fold operating point at 5 ns using corrected RMSD over $(3,5]$ ns. ``Continue'' predicts late pose retention.}
\label{tab:policy-5ns}
\begin{tabular}{lccc}
\toprule
Decision at 5 ns & Late retained & Late non-retained & Total \\
\midrule
Stop & 1 & 25 & 26 \\
Continue & 12 & 14 & 26 \\
\midrule
Total & 13 & 39 & 52 \\
\bottomrule
\end{tabular}
\end{table}

Stopping 26 trajectories at 5 ns and continuing 26 to 100 ns would require $26\times5+26\times100=2{,}730$ simulated ns instead of 5,200 ns. The nominal saving is therefore 2,470 ns, or 47.5\%. This is 230 ns less computation than the evaluated 20 ns policy on the same 52 systems. The 5 ns policy incorrectly stopped 1 of 13 late-retained trajectories and continued 14 of 39 late-non-retained trajectories. These counts are promising but were obtained after checkpoint comparison; they must not be interpreted as a guaranteed future error rate.

\section{Planned Software Tool and Demonstration Workflow}

The locked retrospective framework is intended to become a lightweight, standalone software tool that reads a running simulation's saved trajectory frames directly -- alongside the molecular-dynamics job itself, with no change to the simulation -- and reports a single stop/continue score together with a confidence range, rather than an opaque model prediction. Its planned auditable sequence is:

\begin{enumerate}
    \item Load a topology and the prefix of a running trajectory (its saved coordinate frames, e.g.\ a DCD file, read alongside the still-running simulation).
    \item Check trajectory quality and periodic-box information.
    \item Reconstruct the ligand as one molecule under periodic boundaries and make its position protein-relative.
    \item Fit protein C$\alpha$ atoms to frame 0 and apply that protein-derived transform to ligand heavy atoms.
    \item Calculate corrected ligand RMSD and unweighted ligand-centroid displacement -- the same closed-form measurement defined in Section~\ref{sec:features}, with no machine-learned step at this stage.
    \item For the exploratory early rule, average corrected RMSD over $(3,5]$ ns. For the inherited primary analysis, average both measurements over $(18,20]$ ns.
    \item For the retrospective audit, use the matching checkpoint's held-out fold, training-only scaler, and evaluation model to reproduce its uncalibrated out-of-fold score.
    \item Apply the fixed 0.5 threshold and report \textbf{stop} or \textbf{continue}, together with the same protein-clustered bootstrap confidence range used throughout this study, pocket-retention context, and an explicit same-trajectory warning.
\end{enumerate}

The current audit contains saved out-of-fold scores from 15 protein-held-out evaluation models at each checkpoint, not a final model trained on the full development cohort. A live tool should therefore be presented as an analysis prototype until one frozen all-development 5 ns model is specified and evaluated prospectively. It should lead with corrected RMSD as the primary reported number, use the fitted model output only as a secondary, uncalibrated decision score, show pocket retention as context, and label ERD as exploratory.

\section{Discussion}

The central finding of this report is deliberately narrow: a single, physically defined geometric observable, computed early in a running trajectory, discriminates whether that same trajectory later exhibits geometric pose retention. This claim does not extend to binding, biological activity, or thermodynamic stability, and it is stated at exactly the scope the evidence supports.

That narrowness is the source of its practical value, not a limitation to apologize for. Screening candidate ligands against a panel of resistance-conferring ABC-transporter pockets (Section~\ref{sec:motivation}) requires deciding, for every docked complex, whether it is worth the cost of a complete production trajectory before any more informative and more expensive follow-up -- longer sampling, free-energy estimation, or eventual experimental testing -- is considered. A geometric pose-retention criterion is already the practical first gate in this kind of pipeline: a complex whose ligand has clearly left its docked pocket by the late window is not one that more expensive downstream analysis will rescue, and a complex that is still there is the reason a project keeps its compute allocation on that system rather than reassigning it elsewhere. The contribution of this report is not to justify that gate on chemical or biological grounds -- it is a screening convenience, not a claim about efficacy -- but to show that the same information the gate already uses can be read early, from a correctly processed trajectory, at a small and honestly quantified false-stop cost.

This is also why the coordinate-correction ablation and the same-trajectory-vs-separate-launch boundary (Section~\ref{sec:results}) matter more than they might for a broader claim. A gatekeeping step used to decide which simulations keep running has to be trustworthy under exactly the conditions it will be used in, not merely correlated with outcome in aggregate. Showing that the signal collapses under naive coordinate handling, and that it does not transfer across independently launched realizations of the same complex, is what makes a narrow geometric criterion usable as a real component of a screening pipeline rather than a suggestive but unverified correlation. A broader, less rigorously bounded claim would in fact be less trustworthy for exactly the deployment role proposed here -- the rigor and the narrowness support each other rather than trading off.

\section{Limitations}

\begin{itemize}
    \item The locked cohort contains 52 trajectories from 15 protein groups for one ligand in predicted ABC-transporter systems. It is small, structurally related, and omits the full membrane environment.
    \item The late outcome is complete simulated coverage plus a geometric RMSD threshold. It does not measure binding affinity, biological activity, thermodynamic stability, or experimental success.
    \item Early and late pose are measured with the same corrected-RMSD family. The result establishes temporal persistence within a trajectory, not a separate hidden mechanism.
    \item Observable sets, controls, and alternative checkpoints were compared retrospectively rather than inside a nested selection procedure. The 20 ns checkpoint is historically inherited, while the promising 5 ns checkpoint was identified after examining the checkpoint curve. Neither is established as optimal.
    \item The expanded 73-trajectory model comparison was performed after the locked analysis and is exploratory. Its nested protein-held-out design limits model-selection leakage, but it is not an external or prospective test and should not be used to select a new headline model from the same data. Of the 21 trajectories added beyond the lock, 20 are new pockets or replicas on already-represented proteins and one is from a single newly added protein, so this audit mainly adds power and a stability check rather than testing generalization to a substantially wider protein panel.
    \item The primary 100 ns cohort is being expanded further, toward a target of 80 trajectories. The same-trajectory-vs-separate-launch comparison reported here uses the original locked $n=52$ and has not been re-derived at the larger scale; a planned extension would add roughly 55 pairs already available on disk plus a smaller number recovered from precomputed measurements where the raw trajectory was not retained, but that extension is not part of the results reported in this manuscript.
    \item The fixed 0.5 score threshold is uncalibrated and was not chosen to satisfy a safety or false-stop constraint. No final deployment model, external test set, or prospective shadow run has yet been evaluated.
    \item Only one separately launched trajectory per complex was available. Dynamical independence was checked directly via near-zero potential-energy fluctuation correlation from an identical starting structure (54 of 61 pairs with sufficient overlap; Section~\ref{sec:results}), which is evidence of independence rather than a formal per-run random-seed record. This is not a complete estimate of replica-to-replica reproducibility.
    \item Bootstrap intervals resample protein groups with saved observables or out-of-fold scores held fixed; they do not repeat the full processing, model fitting, and analysis-selection pipeline.
    \item Early stopping is appropriate only for candidate triage. Selective termination can bias equilibrium or kinetic estimates unless the downstream estimator explicitly accounts for the stopping rule.
    \item Pocket retention and ERD are supporting or exploratory measurements. They do not yet establish a validated recovery or failure mechanism.
\end{itemize}

\section{Conclusion}

This project's contribution is a mechanistic finding about ligand-pocket commitment in this cohort, with a practical triage policy built on top of it:

\begin{quote}
\textbf{In this retrospective cohort, a protein-relative, periodic-boundary-consistent ligand-pose signal was already detectable by the 5 ns checkpoint and discriminated late pose retention within the same running trajectory. The same descriptor did not transfer reliably from a separately launched trajectory, indicating that the forecast is trajectory-conditioned rather than an intrinsic stability prediction for the starting complex.}
\end{quote}

Across 52 trajectories from 15 protein groups, corrected RMSD achieved AUROC 0.856 (95\% CI 0.742--0.935) at 5 ns and 0.860 (0.747--0.951) at 20 ns. The corresponding raw AUROCs were 0.835 and 0.825 in the expanded 73-trajectory master cohort. At 20 ns, the PBC-naive and ligand-self-aligned controls achieved 0.550 and 0.507, respectively, while corrected RMSD from a separately launched trajectory achieved 0.546. These results support an early, coordinate-consistent, same-trajectory persistence signal; they do not establish molecular stability or replica-level predictability.

A post-lock model-complexity audit on the expanded 73-trajectory cohort further showed that corrected RMSD retained higher AUROC than nested temporal logistic, spline, random-forest, and gradient-boosting alternatives. This result suggests that physically correct representation contributes more than added classifier complexity in the current dataset and strengthens the case for a transparent persistence rule.

At the exploratory fixed 0.5 single-observable score threshold, stopping 26 trajectories at 5 ns would save 2,470 of 5,200 nominal simulated ns (47.5\%) while stopping 1 of 13 late-retained trajectories. Because the 5 ns checkpoint was selected after examining multiple checkpoints, this operating point is neither prospectively validated, calibrated, nor safety-constrained. The stronger contribution is therefore the evidence that useful same-run persistence information may emerge within 5 ns, together with an explicit cost--error accounting. The next decisive test is a frozen 5 ns prospective evaluation on new trajectories.

\section*{Author Contribution and External Assistance}

\textbf{Complete before submission.} List the named human contributions to structure preparation, docking, molecular-dynamics production, analysis, supervision, and final manuscript approval. Any use of generative-AI tools must be disclosed separately and accurately in the form required by the competition, including the tool, version, purpose, and extent of use; AI-generated prose must not be submitted as the student's independent writing.

\section*{Acknowledgements}

\textbf{Complete before submission.} Acknowledge advisors, collaborators, computing resources, and third-party software.

\clearpage
\begin{thebibliography}{99}

\bibitem{meng2011docking}
X.-Y. Meng, H.-X. Zhang, M. Mezei, and M. Cui,
``Molecular docking: a powerful approach for structure-based drug discovery,''
\emph{Current Computer-Aided Drug Design}, vol. 7, pp. 146--157, 2011.

\bibitem{hollingsworth2018md}
S. A. Hollingsworth and R. O. Dror,
``Molecular dynamics simulation for all,''
\emph{Neuron}, vol. 99, pp. 1129--1143, 2018.

\bibitem{jumper2021alphafold}
J. Jumper et al.,
``Highly accurate protein structure prediction with AlphaFold,''
\emph{Nature}, vol. 596, pp. 583--589, 2021.

\bibitem{mirdita2022colabfold}
M. Mirdita et al.,
``ColabFold: making protein folding accessible to all,''
\emph{Nature Methods}, vol. 19, pp. 679--682, 2022.

\bibitem{krivak2018p2rank}
R. Kriv\'{a}k and D. Hoksza,
``P2Rank: machine learning based tool for rapid and accurate prediction of ligand binding sites from protein structure,''
\emph{Journal of Cheminformatics}, vol. 10, art. 39, 2018.

\bibitem{trott2010vina}
O. Trott and A. J. Olson,
``AutoDock Vina: improving the speed and accuracy of docking with a new scoring function, efficient optimization, and multithreading,''
\emph{Journal of Computational Chemistry}, vol. 31, pp. 455--461, 2010.

\bibitem{kabsch1976}
W. Kabsch,
``A solution for the best rotation to relate two sets of vectors,''
\emph{Acta Crystallographica Section A}, vol. 32, pp. 922--923, 1976.

\bibitem{WHO2022FPPL}
World Health Organization,
\emph{WHO Fungal Priority Pathogens List}, 2022.
Available: \url{https://www.who.int/publications/i/item/9789240060241}

\bibitem{li_azole}
Y. Li, C. Hind, J. Furner-Pardoe, J. M. Sutton, and K. M. Rahman,
``Understanding the mechanisms of resistance to azole antifungals in \emph{Candida} species,''
\emph{JAC-Antimicrobial Resistance}, vol. 7, no. 3, art. dlaf106, 2025.

\bibitem{bolhuis2002tps}
P. G. Bolhuis, D. Chandler, C. Dellago, and P. L. Geissler,
``Transition path sampling: throwing ropes over rough mountain passes, in the dark,''
\emph{Annual Review of Physical Chemistry}, vol. 53, pp. 291--318, 2002.

\bibitem{e2006towards}
W. E and E. Vanden-Eijnden,
``Towards a theory of transition paths,''
\emph{Journal of Statistical Physics}, vol. 123, pp. 503--523, 2006.

\bibitem{peters2016reaction}
B. Peters,
``Reaction coordinates and mechanistic hypothesis tests,''
\emph{Annual Review of Physical Chemistry}, vol. 67, pp. 669--690, 2016.

\bibitem{jing2026embs}
A. Jing,
``A Physics-Informed Machine Learning Framework for Screening Antifungal Resistance-Mediating ABC Transporter Pocket Stability,''
\emph{Proceedings of the 2026 IEEE Engineering in Medicine and Biology Society Conference}, 2026.
\textbf{Publication status to be confirmed before submission.}

\bibitem{eastman2024openmm}
P. Eastman et al.,
``OpenMM 8: Molecular Dynamics Simulation with Machine Learning Potentials,''
\emph{Journal of Physical Chemistry B}, vol. 128, no. 1, pp. 109--116, 2024.

\bibitem{maier2015ff14sb}
J. A. Maier, C. Martinez, K. Kasavajhala, L. Wickstrom, K. E. Hauser, and C. Simmerling,
``ff14SB: Improving the Accuracy of Protein Side Chain and Backbone Parameters from ff99SB,''
\emph{Journal of Chemical Theory and Computation}, vol. 11, no. 8, pp. 3696--3713, 2015.

\bibitem{wang2004gaff}
J. Wang, R. M. Wolf, J. W. Caldwell, P. A. Kollman, and D. A. Case,
``Development and testing of a general amber force field,''
\emph{Journal of Computational Chemistry}, vol. 25, no. 9, pp. 1157--1174, 2004.

\bibitem{jakalian2002am1bcc}
A. Jakalian, D. B. Jack, and C. I. Bayly,
``Fast, efficient generation of high-quality atomic charges. AM1-BCC model: II. Parameterization and validation,''
\emph{Journal of Computational Chemistry}, vol. 23, no. 16, pp. 1623--1641, 2002.

\bibitem{jorgensen1983tip3p}
W. L. Jorgensen, J. Chandrasekhar, J. D. Madura, R. W. Impey, and M. L. Klein,
``Comparison of simple potential functions for simulating liquid water,''
\emph{Journal of Chemical Physics}, vol. 79, no. 2, pp. 926--935, 1983.

\end{thebibliography}

\end{document}
